\section{不变量：证明某件事做不到}
\label{sec:17-Synthesis/02-Recurring-Ideas/04}

一张 \(8\times 8\) 的棋盘，剪掉对角线上的两个角格，剩下 62 格。能不能用 31 张 \(1\times 2\) 的骨牌恰好铺满？动手试，能铺到剩两格空在天南地北；换个铺法，还是剩两格。试上两个小时，会开始怀疑是不是自己不够聪明。

不是。把棋盘按黑白相间涂色，每张骨牌无论怎么放，一定同时盖住一黑一白，所以 31 张骨牌盖住的黑白数必然相等。而对角线上的两个角同色，剪掉之后剩下 32 格一种颜色、30 格另一种。铺不满。

这个论证有一点值得停下来看：它没有考察任何一种具体铺法，却对\emph{所有}铺法给出了结论。三行字排除了一个天文数字规模的搜索空间。做到这一点的东西是一个量——黑格数减白格数——它在每一次放置骨牌的操作下保持不变，而初始值与目标值对不上。

隔壁工位在做一件形式上完全不同的事：验证一个非线性控制器是否全局稳定。系统是六维的，解析解没有，仿真只能覆盖有限多条初始轨迹。工程师写下一个二次型 \(V(x)=x^{\trans}Px\)，验证沿着闭环动力学 \(\dot V(x)<0\) 对所有非零 \(x\) 成立，然后声称所有轨迹都趋于原点——同样没有解方程，同样一举处理了无穷多种情况。

这两件事是同一个动作。

\subsection{一个动作}

第四个反复出现的数学动作是：\emph{与其追踪系统怎样变化，不如找一个在变化中保持不变、或者只朝一个方向变的量。}有了这样一个量，状态空间就被它的等值面分层，而演化被限制在层内或只能朝一侧穿过。凡是与当前层不在同一侧的目标，一律不可达——不需要检查任何一条具体路径。

\begin{center}
\begin{tikzpicture}[x=1cm,y=1cm,font=\small,>=stealth]
  \fill[blue!8] (3.0,2.2) ellipse (2.6 and 1.7);
  \draw[blue!55!black,very thick] (3.0,2.2) ellipse (2.6 and 1.7);
  \draw[blue!35!black,thin] (3.0,2.2) ellipse (1.5 and 1.0);
  \node[font=\footnotesize,blue!50!black,anchor=north] at (3.0,0.42) {\(\{\,V \le c\,\}\)};
  % 内部箭头：V 不增
  \draw[black!55,->] (4.6,2.9) -- (3.9,2.6);
  \draw[black!55,->] (1.5,1.6) -- (2.2,1.9);
  \draw[black!55,->] (3.1,3.6) -- (3.05,3.0);
  \draw[black!55,->] (2.6,0.9) -- (2.75,1.4);
  % 起点与目标
  \fill[green!45!black] (4.2,3.0) circle (0.075);
  \node[font=\footnotesize,green!35!black,anchor=south] at (4.2,3.12) {初始状态};
  \fill[red!75!black] (7.9,3.2) circle (0.075);
  \node[font=\footnotesize,red!70!black,anchor=west] at (8.05,3.2) {想到达的状态};
  \draw[black!45,dashed,->] (4.2,3.0) -- (7.75,3.18);
  \draw[red!75!black,very thick] (4.98,2.9) -- (5.42,3.26);
  \draw[red!75!black,very thick] (4.98,3.26) -- (5.42,2.9);
  \node[font=\footnotesize,black!70,anchor=west,align=left] at (6.3,1.5)
    {\(V\) 沿演化不增\\[2pt] \(\Rightarrow\) 这条线越不过去};
\end{tikzpicture}

\emph{不变量把状态空间切成层。演化被关在自己那一层里，或者只能朝一个方向穿层。所有``到不了''的结论都是这张图：不必看路径长什么样，只要看两端在不在同一侧。}
\end{center}

\subsection{力量来自否定}

构造性的结论和否定性的结论，成本差得很远。要证明某件事\emph{做得到}，给出一个做法就行；要证明某件事\emph{做不到}，得对所有可能的做法负责，而可能的做法通常有无穷多种或指数多种。不变量是少数几种能以常数成本完成后一件事的工具。

\hyperref[sec:01-Mathematical-Language/04-Induction-and-Recursion/04]{``不变量与算法正确性''}把这套方法讲成一件日常的事：循环不变式在进入循环前成立、每轮迭代后仍成立，于是退出时它连同退出条件一起给出正确性。这里的``不变''是相对于一次循环体的执行说的，与棋盘上相对于一次放置骨牌说的完全同构。同一份逻辑还给出终止性：找一个非负整数在每轮严格减少，循环就不可能无限跑下去——这时用的是单调而不是守恒，但它做的是同一件事。

\hyperref[sec:04-Discrete-Mathematics/03-Graph-Theory/07]{``欧拉回路与一笔画''}给的是这类论证最漂亮的一个例子。一笔画走过一条边，就在它的两个端点各消耗一个度；因此除起点终点外每个顶点的度必须是偶数。奇度顶点的个数是那个不变量，它在``已经画了多少''的过程中受到严格约束，于是七桥问题在还没开始搜索时就已经有答案。\hyperref[sec:04-Discrete-Mathematics/03-Graph-Theory/05]{``图着色与冲突安排''}走的是同一条路：色数给出的下界，同样是一个不依赖具体安排方式的量。

优化里的对偶是这件事的连续版本。\hyperref[ch:075]{``对偶：为最优性建立下界证书''}给出的可行对偶解是一张证书——它不告诉你怎么做得更好，它告诉你不可能做得比这个数更好。证书的价值恰恰在于它是否定性的：有了它，搜索可以停下来。

\subsection{同一个动作的其它名字}

Lyapunov 函数是不变量在动力系统里的标准形态，只是把``不变''放宽成``单调不增''。\hyperref[sec:15-Differential-Equations-and-Dynamical-Systems/02-Stability-and-Dynamical-Systems/02]{``稳定性与 Lyapunov 函数''}说明它为什么比线性化判据更强：线性化只在不动点的一个邻域里说话，而一个在整个区域上有效的 \(V\) 直接给出那个区域内的全局结论。上一篇的迭代收敛论证也可以套进这个框架——压缩条件就是以到不动点的距离作为 Lyapunov 函数，并要求它每步至少乘上 \(\gamma\)。

物理里的守恒量与对称性是这个动作最深的一层。\hyperref[sec:12-Real-Analysis-Support/10-Essential-Structures/01]{``对称性怎样产生守恒律''}给出的通道是单向而系统的：作用量在一个连续变换下不变，就有一个对应的守恒量。这是目前唯一一条能\emph{批量生产}不变量的路子，其余情形基本靠对具体问题的洞察。同一条思路在代数化之后就是群作用下的不变函数，\hyperref[sec:04-Discrete-Mathematics/06-Groups-and-Symmetry/04]{``不变函数与等变映射：对称性怎样进入计算''}把``不变''和``等变''分清：前者要求作用之后值不变，后者要求作用与映射可交换。深度学习里把对称性写进架构，用的是后者——\hyperref[sec:14-Deep-Learning/04-Convolutional-Networks/03]{``下采样用分辨率换取不变性''}说明卷积层是平移等变的，而池化把等变逐步换成不变，这个换取有代价：位置信息被主动丢掉了。

概率里对应的东西是鞅。\hyperref[sec:06-Random-Processes/04-Applications/02]{``Martingale：公平游戏与停时''}给出的``条件期望不变''是不变量在随机世界里的正确翻译——单个轨迹当然会变，不变的是关于未来的条件期望。可选停时定理说的是在什么条件下这个不变性能挺过一个随机的停止时刻，而它的条件不是形式主义：加倍下注策略在有限步内几乎必然把亏损扳平，看上去凭空造出了正期望，破绽正在于那个停时不满足定理的条件——它需要无界的本金。同一个不变量在别处叫平稳性，\hyperref[sec:06-Random-Processes/05-Continuous-Time-Markov-Chains-and-Queues/03]{``连续时间平稳分布与流量平衡''}里的流量平衡方程，说的就是``流进等于流出''这个在演化下保持的账。

偏微分方程的适定性论证也是同一套。\hyperref[sec:15-Differential-Equations-and-Dynamical-Systems/03-PDE-Introduction/03]{``适定性与能量方法''}用的是能量泛函：证明它随时间不增，于是解不可能爆炸，而且两个解的差的能量从零出发只能保持为零，唯一性顺手得到。整个论证里没有出现过解的表达式。

最后一处形态不太像不变量，但机制完全一样。神经网络的参数有大量等价对称——隐藏单元可以置换，某些层的尺度可以互相搬移，而函数完全不变。\hyperref[sec:14-Deep-Learning/10-Interpretability-Mathematics/01]{``神经网络的函数空间视角''}由此给出一把很好用的筛子：凡是在这些等价变换下会改变的量，描述的是这一次训练落在了哪个坐标上，而不是模型学到了什么。用它去看``某个神经元代表什么''这类说法，多数会被筛掉。

\subsection{代价：没有通用的找法}

这个动作的力量是真的，它的代价也是真的，而且必须讲明白：\emph{给定一个问题，没有通用方法找到合适的不变量。}

对称性那条通道是例外，它把``找守恒量''换成了``找对称性''，而对称性有时是能看出来的。除此之外，情况并不乐观。Lyapunov 函数有逆定理——渐近稳定的系统必定存在某个 Lyapunov 函数——但那个定理不给构造，实践中仍然靠猜一个二次型然后解线性矩阵不等式，猜不中就换一个形式。组合问题里的染色论证同样如此，黑白染色适用于骨牌，换成 \(1\times 3\) 的骨牌就得换成三色，而``该用几种颜色、怎么涂''没有算法。

所以这个动作在实际工作中的位置是特定的：它不是遇到问题时的第一手，而是在反复尝试构造失败之后应当想起的第二手。``试了很久做不出来''本身就是一条信号，提示应当掉头去找一个不变量，把``做不出来''升级成``做不到''。反过来，找不到不变量不能证明事情做得到——这是一个不对称的工具，只在一个方向上给出结论。

\subsection{边界}

不变量给出的结论强，前提也硬：那个量必须在\emph{每一步操作}下都不变或单调。工程实践中最常见的失守是操作集合悄悄变大——证明时假定系统只能做 A、B 两类动作，上线后加了第三类，不变量不再守恒，而基于它的所有结论一起失效，且不会有任何报错。Lyapunov 分析被一个未建模的执行器饱和推翻，守恒量被一个数值积分器的舍入误差慢慢耗散，都属于这一类。

这一篇处理的是``变化中什么不变''。下一篇处理一个表面上更琐碎、实际上更常出问题的问题：两个操作换个顺序做，结果还一样吗。它是这本书里被使用得最频繁而被交代得最少的一件事。
